%***************************************************************************************************
% preamble/listings.tex - jeden wspólny styl listingów (kod i wyjście terminala).
% Użycie: \begin{lstlisting}[caption={Opis}, label={lst:etykieta}] ... \end{lstlisting}
%         \begin{lstlisting}[style=code, language=Python, ...] ... \end{lstlisting}
%         \lstinputlisting[style=code, language=C, caption={...}, label={...}]{plik.c}
%***************************************************************************************************

% Polskie znaki w listingach (pdfLaTeX + UTF-8); mapowanie zdefiniowane tylko tutaj
\lstset{
  inputencoding=utf8,
  extendedchars=true,
  literate=
    {ą}{{\k{a}}}1 {ć}{{\'c}}1 {ę}{{\k{e}}}1 {ł}{{\l{}}}1 {ń}{{\'n}}1
    {ó}{{\'o}}1 {ś}{{\'s}}1 {ź}{{\'z}}1 {ż}{{\.z}}1
    {Ą}{{\k{A}}}1 {Ć}{{\'C}}1 {Ę}{{\k{E}}}1 {Ł}{{\L{}}}1 {Ń}{{\'N}}1
    {Ó}{{\'O}}1 {Ś}{{\'S}}1 {Ź}{{\'Z}}1 {Ż}{{\.Z}}1
}

% Styl bazowy
\lstdefinestyle{base}{
  basicstyle=\ttfamily\footnotesize,
  columns=fullflexible,
  keepspaces=true,
  frame=single,
  rulecolor=\color{lgray},
  breaklines=true,
  postbreak=\mbox{\textcolor{dgray}{$\hookrightarrow$}\space},
  showstringspaces=false,
  captionpos=t,
  aboveskip=1em,
  belowskip=1em,
}

% Wyjście terminala / sesja konsoli (styl domyślny)
\lstdefinestyle{terminal}{
  style=base,
  backgroundcolor=\color{gray!6},
}

% Kod źródłowy (z numerami wierszy i kolorami)
\lstdefinestyle{code}{
  style=base,
  numbers=left,
  numberstyle=\tiny\color{dgray},
  numbersep=6pt,
  xleftmargin=1.6em,
  framexleftmargin=1.6em,
  keywordstyle=\bfseries\color{Accent},
  commentstyle=\itshape\color{dgray},
  stringstyle=\color{DarkGreen},
}

\lstset{style=terminal}

% Numeracja listingów w obrębie rozdziału (jak rysunki i tabele)
\AtBeginDocument{\counterwithin{lstlisting}{chapter}}
